\chapter{Commutativity}
\label{chap:commutativity}
\label{sec:g-comm}

The pasting step combines several slice measurements $G^x$ into one global
measurement.  This requires that $G^x_g$ and $G^y_h$ approximately commute.
The argument has three steps: the point measurements commute
(Section~\ref{sec:commutativity-points}), the slice measurements commute after
evaluation at random points (Section~\ref{sec:G-commutes-after-evaluation}),
and the Schwartz--Zippel lemma lifts the evaluated commutation to the full
polynomial outcomes (Section~\ref{sec:G-commutes}).

\section{Commutativity of the point measurements}
\label{sec:commutativity-points}

\begin{definition}[Shared diagonal-line sample]\label{def:point-pair-shared-diagonal-line-distribution}
	\lean{MIPStarRE.LDT.CommutativityPoints.pointPairSharedDiagonalLineDistribution}
	\leanok
	\uses{def:low-individual-degree-test}
	Sample independent uniform points $u,v\in\F_q^m$ and a uniform parameter
	$t\in\F_q$.  The associated diagonal-line question is the line with
	direction $v-u$ whose parametrization visits $u$ at $t$ and $v$ at $t+1$.
	When $u=v$ this is the degenerate line $\{u\}$.
\end{definition}

\begin{theorem}[Commutativity of the point measurements]\label{thm:commutativity-points}
	\lean{MIPStarRE.LDT.CommutativityPoints.answerCommutativityPoints, MIPStarRE.LDT.CommutativityPoints.commutativityPointsError}
	\leanok
	\uses{def:good-strategy, def:approx_delta, def:point-pair-shared-diagonal-line-distribution}
	Let $(\psi,A,B,L)$ be an $(\eps,\delta,\gamma)$-good symmetric strategy for
	the $(m,q,d)$ low individual degree test.  On average over independent
	uniform $u,v\in\F_q^m$,
	\[
		(A^u_aA^v_b)\ot I\approx_{32\gamma m}(A^v_bA^u_a)\ot I.
	\]
\end{theorem}
\begin{proof}
	\leanok
	\uses{def:restricted-diagonal-lines-test, prop:simeq-to-approx, prop:cab-approx-delta, prop:triangle-inequality-for-approx_delta}
	The strategy passes the diagonal lines test with probability at least
	$1-\gamma$, hence the $m$-restricted diagonal lines test with probability
	at least $1-\gamma m$.  Thus
	$A^u_a\ot I\simeq_{\gamma m}I\ot L^\ell_{[f(u)=a]}$ on average over
	independent uniform $u,w\in\F_q^m$ with $\ell=\{u+tw\}$, and
	Proposition~\ref{prop:simeq-to-approx} gives
	\[
		A^u_a\ot I\approx_{2\gamma m}I\ot L^\ell_{[f(u)=a]}.
	\]
	For the shared line of
	Definition~\ref{def:point-pair-shared-diagonal-line-distribution}, the
	marginals on $(u,\ell)$ and on $(v,\ell)$ are exactly this distribution.
	Hence
	\begin{align*}
		(A^u_aA^v_b)\ot I
		 & \approx_{2\gamma m}A^u_a\ot L^\ell_{[f(v)=b]}
		\approx_{2\gamma m}I\ot L^\ell_{[f(v)=b]}L^\ell_{[f'(u)=a]} \\
		 & =I\ot L^\ell_{[f'(u)=a]}L^\ell_{[f(v)=b]}
		\approx_{2\gamma m}A^v_b\ot L^\ell_{[f'(u)=a]}
		\approx_{2\gamma m}(A^v_bA^u_a)\ot I,
	\end{align*}
	using projectivity of $L$ in the middle.  Each step also uses
	Proposition~\ref{prop:cab-approx-delta} to insert the relation next to the
	remaining contraction factor, and
	Proposition~\ref{prop:triangle-inequality-for-approx_delta} combines the
	four errors.  In the formalization the line answers are functions on the
	line, so the same argument applies with the relaxed answer alphabet.
\end{proof}

\section{Commutativity of the slice measurements after evaluation}
\label{sec:G-commutes-after-evaluation}

Throughout the rest of this chapter, $m\ge1$, $(\psi,A,B,L)$ is an
$(\eps,\delta,\gamma)$-good symmetric strategy for the $(m+1,q,d)$ low
individual degree test, and $\{G^x\}_{x\in\F_q}\subseteq\polysub mqd$ is a
family of projective sub-measurements with the following three properties.
\begin{enumerate}
	\item (Consistency with~$A$) On average over $(u,x)\sim\F_q^{m+1}$,
	      $A^{u,x}_a\ot I\simeq_\zeta I\ot G^x_{[g(u)=a]}$.
	\item (Strong self-consistency) On average over $x\sim\F_q$,
	      $G^x_g\ot I\approx_\zeta I\ot G^x_g$.
	\item (Boundedness) There are positive semidefinite $Z^x$ with
	      $\E_x\bra\psi(I-G^x)\ot Z^x\ket\psi\le\zeta$ and
	      $Z^x\ge\E_uA^{u,x}_{g(u)}$ for every $x$ and $g\in\polyfunc mqd$.
\end{enumerate}

\begin{definition}[Slice-family hypotheses]\label{def:slice-family-hypotheses}
	\lean{MIPStarRE.LDT.IdxPolyFamily.ConsistentWithPoints, MIPStarRE.LDT.IdxPolyFamily.StronglySelfConsistent, MIPStarRE.LDT.IdxPolyFamily.SliceBoundednessInput}
	\leanok
	\uses{def:simeq, def:strong-self-consistency, def:polymeasurements}
	The three properties above are the consistency, strong self-consistency,
	and boundedness hypotheses on a slice family.
\end{definition}

\begin{lemma}[Boundedness controls rejection]\label{lem:boundedness-rejection}
	\lean{MIPStarRE.LDT.Commutativity.GCommStability.Scalar.scalar_pointwise_cauchy_schwarz_bound, MIPStarRE.LDT.Commutativity.gCommStability_scalar, MIPStarRE.LDT.Commutativity.gCommStabilityTwo_raw_scalar}
	\leanok
	\uses{def:slice-family-hypotheses}
	Let $F$ be a projection and $\{R_i\}$ a sub-measurement on the first
	local space, and let $0\le\overline A_i\le Z$ act on the second.  Then
	$\sum_iFR_iF\ot\overline A_i\le F\ot Z$.  Consequently, if $F^x=I-G^x$,
	the sub-measurements $R^x$ do not depend on the evaluation question $u$,
	and $\overline A_i=\E_uA^{u,x}_{g_i(u)}$, then
	\[
		\Bigl|\E_x\sum_i\bra\psi R_i^xF^x\ot\overline A_i\ket\psi\Bigr|
		\le\sqrt\zeta.
	\]
\end{lemma}
\begin{proof}
	\leanok
	Positivity and $\sum_iR_i\le I$ give
	$\sum_iFR_iF\ot\overline A_i\le F(\sum_iR_i)F\ot Z\le F\ot Z$.  For the
	averaged bound apply Cauchy--Schwarz with factors $\sqrt{R_i}\ot I$ and
	$\sqrt{R_i}F\ot\overline A_i$; the first squared norm is at most one and
	the second is at most $\E_x\bra\psi F^x\ot Z^x\ket\psi\le\zeta$.  The two
	instances used below take $R^y_h=\E_{u,x}\sum_aP_aG^y_hP_a$ and
	$\widetilde R^x_g=\E_{v,y}\sum_bQ_bG^x_gQ_b$ in the notation of the next
	proof.
\end{proof}

\begin{lemma}[Commutativity of~$G$ after evaluation]\label{lem:comm-data-processed-g}
	\lean{MIPStarRE.LDT.Commutativity.commDataProcessedG_of_commutativityPoints, MIPStarRE.LDT.Commutativity.CommDataProcessedGConclusion, MIPStarRE.LDT.Commutativity.commDataProcessedGError}
	\leanok
	\uses{def:slice-family-hypotheses, thm:commutativity-points}
	Let $\nu=48m(\gamma^{1/2}+\zeta^{1/2})$.  On average over independent
	uniform $(u,x),(v,y)\sim\F_q^{m+1}$,
	\[
		G^x_{[g(u)=a]}G^y_{[h(v)=b]}\ot I\approx_\nu
		G^y_{[h(v)=b]}G^x_{[g(u)=a]}\ot I.
	\]
\end{lemma}
\begin{proof}
	\leanok
	\uses{lem:boundedness-rejection, prop:cons-sub-meas, prop:two-notions-of-self-consistency, prop:two-notions-of-self-consistency-after-evaluation, prop:closeness-of-ip, prop:switch-sandwich}
	For fixed $u,v,x,y$ write $P_a=G^x_{[g(u)=a]}$, $Q_b=G^y_{[h(v)=b]}$,
	$P=G^x$, $Q=G^y$, $A_a=A^{u,x}_a$, and $B_b=A^{v,y}_b$, and let
	$\mathcal E(X_{a,b})$ denote the average of
	$\sum_{a,b}\bra\psi X_{a,b}\ket\psi$.  Proposition~\ref{prop:cons-sub-meas}
	gives $P_a\ot I\approx_{4\zeta}P\ot A_a$ and
	$Q_b\ot I\approx_{4\zeta}Q\ot B_b$.
	Propositions~\ref{prop:two-notions-of-self-consistency}
	and~\ref{prop:two-notions-of-self-consistency-after-evaluation} give
	$P_a\ot I\approx_\zeta I\ot P_a$ and the same for $Q_b$.
	Theorem~\ref{thm:commutativity-points} for $m+1$ variables gives
	$I\ot A_aB_b\approx_{32\gamma(m+1)}I\ot B_bA_a$.

	Expanding the commutator and exchanging $(u,x,a)$ with $(v,y,b)$ shows
	that the commutator error equals $2(S-T)$ with
	$S=\mathcal E(Q_bP_aQ_b\ot I)$ and $T=\mathcal E(P_aQ_bP_aQ_b\ot I)$.
	The chain
	\begin{align*}
		T & \approx\mathcal E(P_aQ_bP_aQ\ot B_b)
		\approx\mathcal E(P_aQ_bP_a\ot B_b)
		\approx\mathcal E(P_aQ_bP\ot B_bA_a)
		\approx\mathcal E(P_aQ_b\ot B_bA_a)                            \\
		  & \approx\mathcal E(P_aQ_b\ot A_aB_b)
		\approx\mathcal E(P_aQ_b\ot B_b)
		\approx\mathcal E(P_aQ_b\ot I)\approx S
	\end{align*}
	uses, in order, consistency, the first rejection bound of
	Lemma~\ref{lem:boundedness-rejection}, consistency, the second rejection
	bound, point commutativity, consistency twice, and finally
	Proposition~\ref{prop:switch-sandwich} with self-consistency.  Each step
	is an application of Proposition~\ref{prop:closeness-of-ip}, whose
	normalization hypotheses follow from projectivity and sub-measurement
	normalization.

	The simplified exposition collects the losses as
	$|T-S|\le13\sqrt\zeta+\sqrt{32\gamma(m+1)}$ and obtains the constant
	$26m$.  The formalization follows the same chain with the coarser
	per-step accounting of the original argument, which bounds the point swap
	by $6\sqrt{\gamma(m+1)}$ and gives
	$|T-S|\le12\sqrt\zeta+12\sqrt{\gamma(m+1)}$; since $m+1\le2m^2$ for
	$m\ge1$, the commutator error is at most $48m(\sqrt\gamma+\sqrt\zeta)$.
	Only the downstream bound of Theorem~\ref{thm:com-main} is used later, and
	it is the same for both constants.
\end{proof}

\section{Commutativity of the slice measurements}
\label{sec:G-commutes}

\begin{theorem}[Commutativity of~$G$]\label{thm:com-main}
	\lean{MIPStarRE.LDT.Commutativity.comMain_of_commutativityPoints, MIPStarRE.LDT.Commutativity.ComMainConclusion, MIPStarRE.LDT.Commutativity.comMainError}
	\leanok
	\uses{def:slice-family-hypotheses, lem:comm-data-processed-g}
	Let
	\[
		\nu=30m\bigl(\gamma^{1/4}+\zeta^{1/4}+(d/q)^{1/4}\bigr).
	\]
	On average over independent uniform $(u,x),(v,y)\sim\F_q^{m+1}$,
	\[
		G^x_gG^y_h\ot I\approx_\nu G^y_hG^x_g\ot I.
	\]
\end{theorem}
\begin{proof}
	\leanok
	\uses{cor:schwartz-zippel-individual, prop:closeness-of-ip, prop:switch-sandwich}
	The commutator error is at most $4$ by sub-measurement normalization, so
	the claim is immediate unless $\gamma,\zeta,d/q<1$.  Expanding the
	commutator gives
	\[
		2\,\E_{x,y}\sum_{g,h}\bra\psi G^x_gG^y_hG^x_g\ot I\ket\psi
		-2\,\E_{x,y}\sum_{g,h}\bra\psi G^x_gG^y_hG^x_gG^y_h\ot I\ket\psi.
	\]
	The first term is within $O(\sqrt\zeta)$ of $\bra\psi G\ot G\ket\psi$,
	where $G=\E_xG^x$, by Proposition~\ref{prop:switch-sandwich} and strong
	self-consistency.  For the second term, move one factor $G^x_g$ to the
	second tensor factor by Proposition~\ref{prop:closeness-of-ip} and
	evaluate it at a random point $u$.  Distinct $g\ne g'$ agree at $u$ with
	probability at most $md/q$ by
	Corollary~\ref{cor:schwartz-zippel-individual}, so this costs $md/q$.
	Two more applications of Proposition~\ref{prop:closeness-of-ip} and a
	second Schwartz--Zippel evaluation at an independent point $v$ reach the
	fully evaluated quantity
	\[
		\E_{u,v,x,y}\sum_{a,b}\bra\psi
		G^x_{[g(u)=a]}G^y_{[h(v)=b]}G^x_{[g(u)=a]}\ot G^y_{[h(v)=b]}\ket\psi.
	\]
	Lemma~\ref{lem:comm-data-processed-g} commutes the evaluated factors at
	cost $\sqrt{48m(\sqrt\gamma+\sqrt\zeta)}$, and one more use of
	self-consistency returns to $\bra\psi G\ot G\ket\psi$.  Collecting the
	two Schwartz--Zippel losses, the square-root losses, and the evaluated
	commutation error under $\gamma,\zeta,d/q<1$ gives the stated bound.

	The simplified exposition obtains the same bound by a lifting inequality:
	for evaluation maps with collision probability $\theta<1$ the mirror
	operator $\sum_i(P_i\ot I-I\ot P_i)^2$ decreases by at most the factor
	$1-\theta$ under evaluation, which gives
	$\E_{x,y}C(G^x,G^y)\le8\chi+68\zeta$ with
	$\chi=26m(\sqrt\gamma+\sqrt\zeta)$ when $\theta=md/q\le\frac12$.  The
	formalization uses the two-step Schwartz--Zippel evaluation described
	above.
\end{proof}
